\begin{frame}{왜 궤적을 보나: 누적 보상은 프록시}
  13.1절이 수학으로 미리 말한 바로 그 사실
  ($k_p < 5$ 면 제어기의 ``줄''이 중력 ``역줄''을 이기지 못해 위쪽 균형이
  불안정)을, 여기선 \emph{실험으로} 확인한 셈.
  \vskip0.6em
  \begin{block}{왜 ``궤적을 보는'' 단계가 필요한가}
    \textbf{누적 보상은 진짜 과제의 프록시(proxy)이기 때문}.
    ``누적 벌 최소화''와 ``위쪽을 세운 채 유지''는 \emph{대부분} 같은 방향을
    가리키지만 항상 같은 것은 아님. 프록시를 최적화한 뒤 \emph{진짜 목표}를
    한 번 확인하는 습관이 없으면 ``점수는 좋은데 과제를 못 한'' 정책을
    배포하게 됨 --- Week 12의 \kb{리워드 해킹} 방어선이 1관절 진자에서부터
    시작.
  \end{block}
\end{frame}
